\chapter{Espérance, variance et sommes de variables}
\section*{Exercice 22.1 -- Calcul direct}
[Base] $X$ prend les valeurs $-1,0,2$ avec probabilités $1/4,1/2,1/4$. Calculer son espérance et sa variance.

\section*{Exercice 22.2 -- Transformation affine}
[Base] Sachant $E(X)=3$, $V(X)=4$, calculer l'espérance et l'écart type de $Y=5-2X$.

\section*{Exercice 22.3 -- Somme dépendante}
[Critique] Soit $X$ une Bernoulli de paramètre $1/2$ et $Y=1-X$. Comparer $V(X+Y)$ et $V(X)+V(Y)$.

\section*{Exercice 22.4 -- Binomiale}
[Méthode] Pour $X\sim\mathcal B(100,0{,}3)$, calculer $E(X)$, $V(X)$ et la variance de $X/100$.

\section*{Exercice 22.5 -- Sans indépendance}
[Preuve] Dans une classe fictive de 20 élèves, chacun a probabilité $1/5$ d'oublier un livre. Trouver l'espérance du nombre d'oublis sans supposer l'indépendance.

\section*{Exercice 22.6 -- Moindres carrés}
[Approfondissement] Montrer que $E((X-a)^2)=V(X)+(E(X)-a)^2$. Quel réel $a$ minimise cette quantité ?
